% Independent algebraic audit.  Standard amsmath/amsthm notation only.
% Parent paper supplies theorem/proposition/lemma environments and bibliography.
\subsection{Algebraic conventions and the audited input}
All categories in this subsection are small.  We write
$P_x=\mathbb Z[A(x,-)]\in\mathrm{Ab}^A$ and
$T_i(F)=L_i\operatorname{colim}_A F$.  An ordinary arrow $f:x\to y$
induces the morphism of covariant modules
$P_y\to P_x$, $t\mapsto tf$.  The map $\epsilon_x:P_x\to
\underline{\mathbb Z}$ is the natural augmentation.  For
$L\subseteq P_x$, the notation $\epsilon(L)$ means the image of
$\operatorname{colim}_A L\to\operatorname{colim}_A P_x=\mathbb Z$;
it does not assert that one object realizes every element of this image.
Every $P_x$ is projective and has colimit $\mathbb Z$, by the linear
Yoneda lemma and the initial object $(x,1_x)$ of $x\backslash A$.
The comparison $T_i(\mathbb Z[H])\cong
H_i(B\operatorname{El}_A H;\mathbb Z)$ is the isomorphism of the
normalized bar and nerve chain complexes sending
$(a_0\to\cdots\to a_i,h)$ to the corresponding simplex of the
category of elements.  In particular, totality implies the acyclicity
of these permutation modules, including when $H$ is empty or has many
components.

We read the proofs of TOTEX A.1--A.4 and B.1--B.4, rather than using
their stated conclusions without their maps.  The part of A.1 used here
is precisely the preceding chain-complex identification.  The part of
A.4 used here is that a total category has common successors and
satisfies I and IV.  Its common-successor argument identifies the two
values of the constant two-element action on the union of two disjoint
principal cosieves; Mayer--Vietoris gives nonzero first homology of the
resulting double.  Applying this argument in the category of elements of
$A(p,-)$ completes a span.  Its pair argument applies the two short
exact sequences below to the quotient by one family of binomial
relations.  These are the exact steps required here.  We also opened
\cite{IsbellMitchell1976}, Lemma 2.2, Proposition 2.3, Theorem 3.1 and
Proposition 3.3.  Lemma 2.2 reduces a finite zigzag in an indecomposable
action by successive commuting squares.  Proposition 2.3, for IV, then
follows by equalizing the two arrows in the resulting cospan.  Theorem
3.1 identifies exactness with componentwise filtered affinization, and
Proposition 3.3 reduces its pair condition for affine arrows to a finite
family of ordinary arrows.  These are the versions, with covariant
actions, used below.

\begin{lemma}[Proved: the row calculation used below]\label{B:row-audit}
Let $a_0,\ldots,a_k:p\to q$ be a finite family of ordinary arrows in
an arbitrary small category, and let $z_i=a_i-a_0$ for $1\leq i\leq k$.
Put
\[
 d:P_q^k\longrightarrow P_p,\qquad
 (t_i)\longmapsto\sum_i t_i z_i,\qquad K=\ker d,
\]
and let $H$ be $A(p,-)$ modulo all relations $ca_i\sim ca_0$.
Then $\mathbb Z[H]=\operatorname{coker}d$ canonically, and
\begin{align*}
 T_1(\mathbb Z[H])&\cong
 \operatorname{coker}(\operatorname{colim}K\xrightarrow{\epsilon^k}\mathbb Z^k),\\
 T_2(\mathbb Z[H])&\cong
 \ker(\operatorname{colim}K\xrightarrow{\epsilon^k}\mathbb Z^k),\\
 T_j(\mathbb Z[H])&\cong T_{j-2}(K)\quad(j\geq3).
\end{align*}
Consequently, for total $A$, $\epsilon^k$ is an isomorphism and $K$
is colimit-acyclic.  For $k=1$ the module
$J_i=P_qz_i\subseteq P_p$ is colimit-acyclic and has zero colimit.
\end{lemma}
\begin{proof}
At each object, the kernel of the map from the free group on a set to
the free group on its quotient is generated by the differences defining
the equivalence relation; a path of elementary relations expresses each
other difference as their sum.  This identifies the displayed cokernel
without any injectivity assumption on the action maps.  Set
$J=\operatorname{im}d$.  Apply the long exact derived-colimit sequences
to $0\to K\to P_q^k\to J\to0$ and
$0\to J\to P_p\to\mathbb Z[H]\to0$.  The map
$\operatorname{colim}J\to\mathbb Z$ is zero, because its composite
with the epimorphism $\mathbb Z^k\to\operatorname{colim}J$ is
$\operatorname{colim}d=0$.  Exactness and the acyclicity of the
representables give all three formulas.  For total $A$ the left sides
vanish in positive degrees by the bar identification.  In the one-row
case, the same two exact sequences give
$\operatorname{colim}J_i=0$ and $T_j(J_i)=0$ for $j>0$.
\end{proof}

\begin{proposition}[Proved: Lead 1]\label{B:lead1}
Let $A$ be total and $a_0,a_1,a_2:p\to q$.  Put
$K_i=\ker(P_q\xrightarrow{z_i}P_p)$ and $L=K_1\cap K_2$.
Then the surjection in TOTEX B.1 is a canonical isomorphism
\[
 T_1\bigl(P_q/(K_1+K_2)\bigr)\ \cong\
 \mathbb Z/\epsilon(L).
\]
\end{proposition}
\begin{proof}
Lemma~\ref{B:row-audit} with one relation gives canonical augmentation
isomorphisms $\operatorname{colim}K_i\cong\mathbb Z$.
Let $S=K_1+K_2\subseteq P_q$.  The exact sequence
\[
 0\longrightarrow L\xrightarrow{\ell\mapsto(\ell,-\ell)}
 K_1\oplus K_2\xrightarrow{(u,v)\mapsto u+v}S\longrightarrow0
\]
identifies $\operatorname{colim}S$ with
$\mathbb Z^2/\epsilon(L)(1,-1)$.  The augmentation to
$\operatorname{colim}P_q=\mathbb Z$ is induced by $(m,n)\mapsto m+n$.
The sequence $0\to S\to P_q\to P_q/S\to0$ identifies its kernel
with $T_1(P_q/S)$.  Sending the class of $(m,-m)$ to
$m\bmod\epsilon(L)$ gives the stated isomorphism.  This is exactly
the map defining the prior surjection, so the result proves its
injectivity, rather than giving an unrelated abstract isomorphism.

For completeness, the second derivation in the prompt also type-checks,
with its necessary sign made explicit.  Put
$J_i=P_qz_i\subseteq P_p$ and let
$D=\{(tz_1,tz_2):t\in P_q\}\subseteq J_1\oplus J_2$.
The map
\[
 P_q\longrightarrow (J_1\oplus J_2)/D,
 \qquad t\longmapsto[(tz_1,0)]
\]
is onto: $[(uz_1,vz_2)]=[((u-v)z_1,0)]$.
Its kernel is $K_1+K_2$, because $(tz_1,0)=(wz_1,wz_2)$
means $w\in K_2$ and $t-w\in K_1$.  Also $D\cong P_q/L$.
Both $J_i$ have zero colimit and are acyclic by
Lemma~\ref{B:row-audit}.  The exact sequence
$0\to D\to J_1\oplus J_2\to (J_1\oplus J_2)/D\to0$
therefore gives $T_1(P_q/S)\cong\operatorname{colim}D
=\mathbb Z/\epsilon(L)$.
With the usual connecting-map convention, this second description
using $t\mapsto[(tz_1,0)]$ is the negative of the first augmentation
description: a lift $u_1+u_2\in S$ maps to the diagonal element
represented by $u_2$, whose augmentation is $-\epsilon(u_1)$.
Using $t\mapsto[(-tz_1,0)]$ makes the two specified isomorphisms agree.
\end{proof}

\begin{proposition}[Proved: Lead 2]\label{B:lead2}
Let $A$ satisfy I; neither connectedness nor nonemptiness is required
in this proposition.  The following are equivalent:
\begin{enumerate}
\item $\operatorname{colim}_A:\mathrm{Ab}^A\to\mathrm{Ab}$ is exact;
\item for every $p,r\in A$ and every $z\in\mathbb Z[A(p,r)]$
with $\operatorname{aug}z=0$, some affine $t:r\to s$ satisfies $tz=0$.
\end{enumerate}
For every small $A$, every such $p,r,z$, and
$L_z=\ker(P_r\xrightarrow{t\mapsto tz}P_p)$, one has canonically
\[
 T_1(P_p/P_rz)\cong\mathbb Z/\epsilon(L_z).
\]
\end{proposition}
\begin{proof}
Set $J=P_rz$.  The morphism $P_r\to P_p$ has zero map on colimits.
The sequences $0\to L_z\to P_r\to J\to0$ and
$0\to J\to P_p\to P_p/J\to0$ therefore give
$T_1(P_p/J)\cong\operatorname{colim}J
=\mathbb Z/\epsilon(L_z)$, just as in
Lemma~\ref{B:row-audit}.  This formula does not use I.

If colimits are exact, the displayed quotient is zero.  I makes the
objects in each connected component have common successors: successively
complete the spans along a finite connecting zigzag.  All objects at
which $L_z$ is nonzero are in the component of $r$.  A finite expression
for $1\in\epsilon(L_z)$ can therefore be transported to one object
and added there, producing an affine annihilator.

Conversely, if $\alpha,\beta:p\to r$ are affine arrows, their
difference is such a $z$, so (2) equalizes every parallel affine pair.
Each component of $A$ has common successors by I, and hence so does its
affinization.  The pair condition makes every component of
$\operatorname{Aff}_{\mathbb Z}(A)$ filtered.  To verify the implication
to exactness directly, extend an abelian diagram affinely.  Its ordinary
and affine colimits are canonically equal, since for
$t=\sum n_if_i$ of augmentation one,
$x-tx=\sum n_i(x-f_ix)$ in the direct-sum presentation of the colimit.
Each affine component has exact filtered colimits, and the colimit over
all components is their direct sum, an exact functor in abelian groups.
The empty category gives the zero functor and satisfies both conditions
vacuously.
\end{proof}

\begin{proposition}[Proved: Lead 3]\label{B:lead3}
For every small category $A$, every integer $k\geq0$ and every family
$a_0,\ldots,a_k:p\to q$, put
\[
 d^{\mathrm t}:P_q\longrightarrow P_p^k,
 \quad t\longmapsto(tz_i)_{i=1}^k,
 \qquad Q_a=\operatorname{coker}d^{\mathrm t},
 \quad L=\ker d^{\mathrm t}.
\]
Then canonically
\begin{align*}
 T_1(Q_a)&\cong\mathbb Z/\epsilon(L),\\
 T_2(Q_a)&\cong\ker(\operatorname{colim}L\to\mathbb Z),\\
 T_i(Q_a)&\cong T_{i-2}(L)\qquad(i\geq3).
\end{align*}
If objects in the component of $q$ have common successors, the family
has an affine equalizer if and only if $T_1(Q_a)=0$.
In the covariant-module presentation, $Q_a$ has $k$ generators $x_i$
at $p$ and one relation $\sum_i(a_i-a_0)x_i=0$ at $q$.
\end{proposition}
\begin{proof}
Let $D=\operatorname{im}d^{\mathrm t}$.  In the sequence
$0\to D\to P_p^k\to Q_a\to0$, the map on colimits from $D$
is zero: precompose it with the epimorphism
$\mathbb Z=\operatorname{colim}P_q\to\operatorname{colim}D$,
and obtain the zero column $\operatorname{colim}d^{\mathrm t}$.
Thus $T_1(Q_a)=\operatorname{colim}D$ and
$T_i(Q_a)=T_{i-1}(D)$ for $i\geq2$.
The exact sequence $0\to L\to P_q\to D\to0$ gives the three
asserted formulas.  The common-successor argument in the preceding
proof converts augmentation image $\mathbb Z$ into one affine
equalizer; the reverse implication is immediate.  The presentation
follows by sending the generator $1_q\in P_q(q)$ to the indicated
vector in $P_p(q)^k$.  For $k=0$, $Q_a=0$ and $L=P_q$, so every
formula still holds and $1_q$ is the equalizer of the one-arrow family.
The transposed row $P_q^k\to P_p$ has the permutation cokernel of
Lemma~\ref{B:row-audit}; no identification of these two cokernels is
being used when $k>1$.
\end{proof}

\begin{proposition}[Proved: the precise assertions of Lead 4]\label{B:lead4}
Suppose $A$ is nonempty and connected and satisfies I and IV.  For
$M:A\to\mathrm{Ab}$ and $\alpha,\beta\in M(q)$, equality of their
classes in $\operatorname{colim}^{\mathrm{Set}}_A UM$ implies that
an affine $t:q\to s$ satisfies $t\alpha=t\beta$ in $M(s)$.
Moreover $A$ is exact if and only if, for every such $M,q,\alpha,\beta$,
equality in $\operatorname{colim}^{\mathrm{Ab}}_A M$ implies the
same affine equalizability.
\end{proposition}
\begin{proof}
The first equality puts the two elements in one connected component
$X\subseteq UM$.  We verify the part of
\cite{IsbellMitchell1976}, Lemma 2.2 and Proposition 2.3, being used.
Condition I reduces a zigzag in $X$ to a cospan: at each reduction, a
span of ordinary arrows is replaced by a commuting square.  Thus some
$b,b':q\to r$ satisfy $b\alpha=b'\beta$ in $X(r)$.
Condition IV gives an affine $c:r\to s$ with $cb=cb'$ in
$\mathbb Z[A(q,s)]$.  Put $t=cb=cb'$.  Linearity of $M$ yields
$t\alpha=c(b\alpha)=c(b'\beta)=t\beta$.

If $A$ is exact, Theorem 3.1 of the opened paper makes its affinization
filtered.  The equality of ordinary and affine abelian colimits was
proved in Proposition~\ref{B:lead2}.  In a filtered colimit of groups,
two elements at the same object become equal under one transition,
which here is an affine arrow.  Conversely apply the asserted property
to $M=P_p$ and to any two affine arrows $\alpha,\beta:p\to q$.
Their classes in $\operatorname{colim}P_p=\mathbb Z$ are both $1$,
so the affine category equalizes every parallel pair.  Common
successors and Proposition~\ref{B:lead2}'s argument prove exactness.
\end{proof}

\begin{remark}[Proved: scope of the final sentence of Lead 4]\label{B:lead4-scope}
The comparison map
\[
 \operatorname{colim}^{\mathrm{Set}}_A UM\longrightarrow
 U\operatorname{colim}^{\mathrm{Ab}}_A M
\]
is surjective under the common-successor hypothesis: transport the
finite sum representing an abelian colimit element to one object and
add there.  The two assertions above precisely locate the distinction
in its fibres: set equality is already sufficient, while abelian
equality may identify additional components.  This does not identify
exactness with finite-product preservation of the set-colimit functor.
For example, for $A=\Delta_{\mathrm{inj}}$, $p=[0]$, and $M=P_p$,
the elements $2[a]-[b]$ for distinct $a,b:p\to q$ and $[a]$ have the
same abelian-colimit class.  Ordinary increasing injections preserve
the number and coefficients of nonzero basis terms, so these two
elements lie in different set-colimit components.  Nevertheless $A$
is exact, by the already audited Isbell example.  Thus the final phrase
of Lead 4 is valid as motivation for the difference between the two
colimits, and would be false if read as the asserted equivalence with
siftedness or with injectivity of this comparison.
\end{remark}

\begin{lemma}[Proved: a strengthened normal form for zero classes]
\label{B:zero-normal-form}
Let $A$ be a nonempty category in which every pair of objects has a
common successor, and let $F:A\to\mathrm{Ab}$.  For $\kappa\in F(s)$,
its colimit class is zero if and only if there are an ordinary arrow
$u:s\to y$, finitely many ordinary parallel pairs
$e_j,e'_j:x_j\rightrightarrows y$, and $k_j\in F(x_j)$, such that
\[
 F(u)\kappa=\sum_j\bigl(F(e_j)-F(e'_j)\bigr)k_j\quad\hbox{in }F(y).
\]
In particular Lead 5(ii) holds; its condition I can be weakened to
common successors.
\end{lemma}
\begin{proof}
Define the subgroup
\[
 N_y=\sum_{x\in A}\sum_{e,e'\in A(x,y)}
       \operatorname{im}\bigl(F(e)-F(e')\bigr)\subseteq F(y).
\]
Postcomposition sends $N_y$ into $N_z$, so $N$ is a subfunctor.
Any parallel arrows $y\rightrightarrows z$ induce the same map
$F(y)/N_y\to F(z)/N_z$, by the definition of $N_z$.
Consequently $F/N$ factors through the reachability preorder $D$:
its objects are those of $A$, and it has exactly one arrow $y\to z$
when $A(y,z)\ne\varnothing$.  Common successors make $D$ filtered.
Every element of every $N_y$ dies in $\operatorname{colim}_A F$,
and the universal properties therefore identify
\[
 \operatorname{colim}_A F\cong
 \operatorname{colim}_A(F/N)\cong
 \operatorname{colim}_D\overline F
\]
canonically.  If $\kappa$ dies, its image in $\overline F(s)$ becomes
zero at one object $y$ of the filtered diagram $\overline F$.
Choose any arrow $u:s\to y$ in $A$.  Then $F(u)\kappa\in N_y$,
and membership in a sum of subgroups supplies the finite expression
claimed; coefficients can be absorbed into $k_j$.  Conversely every
summand on the right has zero colimit class, and an arrow preserves
that class.  The empty sum is allowed.
\end{proof}

\begin{proposition}[Proved: Lead 5(i), (iii), and (iv)]\label{B:lead5}
Let $A$ be total and let $a_0,a_1,a_2:p\to q$.
Use $z_i,K_i,K,L$ as in Lemma~\ref{B:row-audit} and
Proposition~\ref{B:lead1}, and put
$K_3=\ker(P_q\xrightarrow{a_2-a_1}P_p)$.
Choose affine pair equalizers $\rho_1\in K_1(s)$,
$\rho_2\in K_2(s)$ and $\rho_3\in K_3(s)$ at a common object $s$.
Then
\[
 \kappa=(\rho_1-\rho_3,\rho_3-\rho_2)\in K(s)
\]
has zero class in $\operatorname{colim}K$.
If an affine $t:s\to r$ kills $\kappa$, then
$t\rho_1=t\rho_2=t\rho_3$ is an affine equalizer of the triple.
Such a $t$ exists whenever the finitely many pairs in a normal form of
Lemma~\ref{B:zero-normal-form} have a common affine equalizer.
In particular, a normal form with at most one summand implies the
existence of such a $t$.
Finally the natural morphism
\[
 \psi:K_1\oplus K_2\oplus K_3\longrightarrow K,
 \qquad(k_1,k_2,k_3)\longmapsto(k_1-k_3,k_2+k_3)
\]
has kernel canonically isomorphic to $L$, and
\[
 T_1(K/\operatorname{im}\psi)\cong\mathbb Z/\epsilon(L).
\]
\end{proposition}
\begin{proof}
The row applied to $\kappa$ is
$\rho_1z_1-\rho_2z_2+\rho_3(z_2-z_1)=0$.
Both coordinates of $\kappa$ have augmentation zero, and
Lemma~\ref{B:row-audit} identifies $\operatorname{colim}K$ with
$\mathbb Z^2$ by those augmentations.  Hence its colimit class
vanishes, proving (i).
If $t\kappa=0$, equality of its first and second coordinates gives
$t\rho_1=t\rho_3=t\rho_2$; the common affine arrow belongs to
$K_1(r)\cap K_2(r)$, proving the first assertion in (iii).
For a normal form $u\kappa=\sum(e_j-e'_j)k_j$, a common affine
equalizer $v$ of the pairs gives $(vu)\kappa=0$.  Condition IV gives
such a $v$ for one pair.  No assertion is made that IV alone
equalizes a finite list of pairs.

The displayed formula for $\psi$ has image in $K$, because
$-k_3z_1+k_3z_2=k_3(a_2-a_1)=0$.  Its kernel consists precisely of
$(\ell,-\ell,\ell)$ with $\ell\in L$; membership in $K_3$
follows from membership in $K_1\cap K_2$.  Write
$S=K_1\oplus K_2\oplus K_3$, $I=\operatorname{im}\psi$ and
$C=K/I$.  Each $K_i$ is colimit-acyclic with augmentation colimit
$\mathbb Z$, by the one-relation case of
Lemma~\ref{B:row-audit}, and $K$ is colimit-acyclic with colimit
$\mathbb Z^2$.  Right exactness applied to $0\to L\to S\to I\to0$
gives
\[
 \operatorname{colim}I
 =\mathbb Z^3/\epsilon(L)(1,-1,1).
\]
The map to $\operatorname{colim}K=\mathbb Z^2$ is induced by
$(a,b,c)\mapsto(a-c,b+c)$.  Its kernel is
$\mathbb Z(1,-1,1)/\epsilon(L)(1,-1,1)$.
The exact sequence $0\to I\to K\to C\to0$ identifies this
kernel with $T_1(C)$, because $T_1(K)=0$.

For the graph interpretation, at each $x$ take a bipartite multigraph
with vertex set $A(q,x)\amalg A(p,x)$ and, for each $g:q\to x$,
three labelled edges $E_{g,i}:g\to ga_i$, $i=0,1,2$.
Distinct edges remain distinct even when their endpoints coincide.
The homomorphism from $K(x)$ to the integral graph cycle group is
\[
 (t_1,t_2)\longmapsto
 \sum_g\left(-(t_{1,g}+t_{2,g})E_{g,0}
             +t_{1,g}E_{g,1}+t_{2,g}E_{g,2}\right).
\]
The boundary at a vertex $g$ is zero, and the boundaries at the
other vertices are exactly $t_1z_1+t_2z_2$.  Conversely the boundary
equation at each $g$ expresses its edge-zero coefficient in terms of
the other two, proving this is an isomorphism.  It is natural under
postcomposition, including maps that identify vertices or edges.
The three two-colour cycle groups give respectively
$(k_1,0)$, $(0,k_2)$ and $(-k_3,k_3)$, whose sum is exactly
$\operatorname{im}\psi$.  Thus the graph description is literal over
$\mathbb Z$, not just an analogy over a field.
\end{proof}

\begin{corollary}[Proved: what the mixed-cycle obstruction says]
\label{B:mixed-cycle-obstruction}
Under the hypotheses of Proposition~\ref{B:lead5}, suppose the triple
has no affine equalizer.  Its class $[\kappa]$ in
$\operatorname{colim}I$ is the nonzero generator of
$\ker(\operatorname{colim}I\to\operatorname{colim}K)
\cong\mathbb Z/\epsilon(L)$.
It cannot have a normal form in which every $k_j$ belongs to $I(x_j)$,
and it cannot have a normal form whose pairs have a common affine
equalizer.  In particular every one-term normal form is impossible.
\end{corollary}
\begin{proof}
The element $\kappa$ is $\psi(\rho_1,-\rho_2,\rho_3)$; its three
augmentations are $(1,-1,1)$.  This gives the asserted generator.
The absence of an affine equalizer means
$\epsilon(L)\ne\mathbb Z$, using common successors, so that
generator is nonzero.  A normal form with coefficients in $I$ would
kill its class already in $\operatorname{colim}I$, a contradiction.
The second assertion follows from Proposition~\ref{B:lead5}(iii).
Thus ``carried around a 1-cycle'' must mean a bar chain with
coefficients in $K$, or equivalently the connecting class in
$T_1(K/I)$, rather than a nontrivial ordinary homology class of $BA$:
the latter is contractible for a total category.
\end{proof}
